\documentclass[aspectratio=169]{beamer}

\usetheme{Madrid}
\usecolortheme{default}
\usepackage[utf8]{inputenc}
\usepackage[T1]{fontenc}
\usepackage[french]{babel}
\usepackage{listings}
\usepackage{xcolor}
\usepackage{amssymb}

\definecolor{codebg}{RGB}{245,245,245}
\lstset{
  basicstyle=\ttfamily\scriptsize,
  backgroundcolor=\color{codebg},
  breaklines=true,
  frame=single,
  columns=fullflexible
}

\title{Intégration de Claude Code dans ARIES}
\subtitle{Détail du travail réalisé --- point d'étape}
\author{Negou Donald}
\institute{Stage --- ARIES (Agent Runtime \& Infrastructure Experimentation System)}
\date{\today}

\begin{document}

\frame{\titlepage}

\begin{frame}{Plan}
\tableofcontents
\end{frame}

% ============================================================
\section{Rappel en une slide}

\begin{frame}{Objectif (rappel, une seule slide)}
Ajouter \textbf{Claude Code} comme troisième harness dans ARIES (aux côtés
d'OpenClaw et Hermes), pour pouvoir le faire tourner sur des tâches
Terminal-Bench 2 et comparer sa télémétrie (latence, CPU, mémoire) aux deux
autres.

\vspace{0.4cm}
Tout ce qui suit est \textbf{ce qui a été fait concrètement} pour y arriver :
code écrit, tests réels exécutés, bugs trouvés et corrigés, état actuel.
\end{frame}

% ============================================================
\section{Code livré}

\begin{frame}[fragile]{Fichiers créés}
\begin{lstlisting}
pkg/harness/claudecode/
    harness.go     Start / Run / Stop (interface AgentHarness)
    config.go      rendu settings.json, wrapper bash, exec Docker

pkg/bridge/claudecodessh/
    bridge.go      serveur SSH, dispatch des commandes
    grammar.go     decodage/classification des commandes bash reelles

cmd/aries-claudecode-probe/          outil de diagnostic (connectivite)
cmd/aries-claudecode-harness-probe/  outil de diagnostic (harness complet)
\end{lstlisting}
Plus modifications de \texttt{cmd/aries/wiring.go} (paire
\texttt{claude-code}/\texttt{claude-code-ssh}) et \texttt{pkg/config/config.go}
(backend \texttt{anthropic}).

\vspace{0.2cm}
\textbf{Tout compile, passe \texttt{go vet}} --- mais chaque étape ci-après a
été \textbf{testée contre de vraies infrastructures} (Docker, SSH, Claude Code
réellement installé), pas seulement relu.
\end{frame}

% ============================================================
\section{Étape 1 --- Prototype de connectivité}

\begin{frame}[fragile]{Étape 1 : valider le mécanisme du pont, isolément}
\textbf{Méthode} (\texttt{cmd/aries-claudecode-probe}) :
\begin{enumerate}
  \item Démarrage d'un vrai sandbox Docker (\texttt{ubuntu:24.04})
  \item Démarrage du pont \texttt{claudecodessh} dessus
  \item Connexion avec un vrai client SSH (\texttt{golang.org/x/crypto/ssh})
  \item Envoi de plusieurs commandes séquentielles sur un même canal
\end{enumerate}

\vspace{0.3cm}
\textbf{Résultat} :
\begin{lstlisting}
PROBE_MARKER_START
hello-from-sandbox
/root
PROBE_MARKER_END
=== RESULT: OK ===
\end{lstlisting}
Le mécanisme de proxy vers le sandbox fonctionne. Outil conservé pour la suite.
\end{frame}

% ============================================================
\section{Étape 2 --- Capture de la grammaire réelle}

\begin{frame}{Étape 2 : observer ce que Claude Code envoie vraiment}
\textbf{Méthode} : remplacer \texttt{/bin/bash} par un script espion (journalise
puis délègue au vrai bash) dans un conteneur jetable avec \textbf{Claude Code
réellement installé et authentifié}. Tâche donnée : ``liste les fichiers, puis
affiche la date''.

\vspace{0.3cm}
\textbf{Hypothèse de départ} : un shell interactif persistant pour toute la
session.

\textbf{Réalité observée} : \textbf{5 invocations séparées} de bash, jamais de
session persistante.
\end{frame}

\begin{frame}[fragile]{Étape 2 : les 5 invocations observées}
\begin{lstlisting}
1) bash -c env
   (sonde initiale)

2) bash -c -l "SNAPSHOT_FILE=.../snapshot-bash-<id>.sh
   source ~/.bashrc ...
   (capture fonctions/alias/PATH, une seule fois par session)

3) bash -c -l "shopt -u extglob ... &&
   eval 'ls -la' < /dev/null &&
   pwd -P >| /tmp/claude-<id>-cwd"
   (par appel d'outil)

4) meme forme pour 'date'
\end{lstlisting}
\textbf{Conclusion} : Claude Code simule la persistance d'état en
\textbf{re-sourçant un fichier snapshot} à chaque appel --- exactement le
modèle d'exécution discrète de Hermes/OpenClaw, pas une session continue.
Le pont a été \textbf{entièrement réécrit} pour refléter cela (plus simple
que la conception initiale).
\end{frame}

% ============================================================
\section{Étape 3 --- Authentification}

\begin{frame}[fragile]{Étape 3 : trouver le bon mécanisme d'authentification}
Trois tentatives, testées une par une :

\vspace{0.2cm}
\textbf{1) Variable d'environnement seule} --- échec
\begin{lstlisting}
export ANTHROPIC_API_KEY=...
claude -p "..." -> "Not logged in - Please run /login"
\end{lstlisting}

\textbf{2) apiKeyHelper lisant la variable d'environnement} --- échec
\begin{lstlisting}
"apiKeyHelper": "echo $ANTHROPIC_API_KEY"
-> "apiKeyHelper failed: did not return a value"
\end{lstlisting}
(Claude Code retire la variable avant d'invoquer le helper, pour éviter une
dépendance circulaire.)

\textbf{3) apiKeyHelper lisant un fichier} --- \textbf{fonctionne}
\begin{lstlisting}
"apiKeyHelper": "cat /chemin/prive/vers/la/cle"
-> authentification reussie
\end{lstlisting}
\end{frame}

% ============================================================
\section{Étape 4 --- Harness complet, de bout en bout}

\begin{frame}{Étape 4 : construire et tester le vrai harness}
\textbf{Image Docker construite} : \texttt{aries-claudecode:test-1}
\begin{itemize}
  \item Base \texttt{node:22-slim} + \texttt{openssh-client}
  \item \texttt{npm install -g @anthropic-ai/claude-code}
  \item Utilisateur non-root fixe, UID 10000 (\texttt{aries})
\end{itemize}

\vspace{0.3cm}
\textbf{Test} (\texttt{cmd/aries-claudecode-harness-probe}) : vrai sandbox +
vrai pont + vrai \texttt{Manager.Start/Run/Stop} + vraie clé API.

\vspace{0.3cm}
\textbf{Premier résultat} : \texttt{"Credit balance is too low"} --- le
mécanisme d'authentification fonctionnait, il manquait du crédit sur le compte
API. Une fois le compte crédité : les bugs suivants sont apparus, un par un.
\end{frame}

% ============================================================
\section{Bugs trouvés et corrigés}

\begin{frame}[fragile]{Bug 1 --- Mauvais nom de champ dans settings.json}
\textbf{Avant} (mécanisme inventé, silencieusement ignoré par Claude Code) :
\begin{lstlisting}
{"env": {"ANTHROPIC_API_KEY_HELPER": "cat ..."}}
\end{lstlisting}
\textbf{Après} (mécanisme réel, confirmé à l'étape 3) :
\begin{lstlisting}
{"apiKeyHelper": "cat ..."}
\end{lstlisting}
Sans erreur explicite --- juste un retour silencieux à \texttt{"Not logged in"}.
\end{frame}

\begin{frame}{Bug 2 --- Permissions non-interactives et propriété des fichiers}
\begin{itemize}
  \item \texttt{--dangerously-skip-permissions} manquant (obligatoire en non
        interactif) --- mais \textbf{refusé sous root} par Claude Code
  \item Les fichiers privés stagés (clé API, identité SSH, mode 0600)
        appartenaient à \texttt{root} par défaut --- illisibles par un
        utilisateur non-root
\end{itemize}

\vspace{0.3cm}
\textbf{Correction} : UID/GID fixes (\texttt{runtimeUID = 10000}), appliqués :
\begin{itemize}
  \item au conteneur (\texttt{container.Config.User})
  \item à chaque fichier de l'archive stagée (\texttt{tar.Header.Uid/Gid})
\end{itemize}
\end{frame}

\begin{frame}[fragile]{Bug 3 --- Arguments du wrapper tronqués}
Le wrapper ne lisait que \texttt{\$1}/\texttt{\$2}, en supposant 2 arguments
(\texttt{flag}, \texttt{script}).

\vspace{0.2cm}
Mais l'invocation réelle du générateur de snapshot en a \textbf{trois} :
\texttt{bash -c -l "<script>"}.

\vspace{0.2cm}
\textbf{Résultat avant correction} : le 3\textsuperscript{e} argument (le vrai script)
était silencieusement perdu --- seul \texttt{"bash -c -l"} arrivait au pont,
rejeté à chaque fois.

\vspace{0.3cm}
\textbf{Correction} : le wrapper traite maintenant \textbf{tous} les arguments
reçus, quel que soit leur nombre.
\end{frame}

\begin{frame}[fragile]{Bug 4 --- Reconstruction ambiguë de la commande}
Même après le bug 3 corrigé : joindre les arguments avec de simples
\textbf{espaces} reste ambigu à re-découper côté pont --- un script réel peut
lui-même commencer par un mot qui ressemble à un drapeau
(\texttt{"-l shopt -u extglob ..."} confondu avec un drapeau \texttt{-l}
suivi d'un script commençant par \texttt{shopt}).

\vspace{0.3cm}
\textbf{Correction} : séparateur non imprimable exact
(\texttt{\textbackslash x1f}, ASCII Unit Separator --- même convention que les
marqueurs de fin d'exécution d'Hermes). Le wrapper joint tous les arguments
originaux avec ce séparateur ; le pont les resépare avec
\texttt{strings.Split}, sans aucune ambiguïté possible.
\end{frame}

\begin{frame}[fragile]{Bug 5 --- Motif de grammaire trop strict}
Le motif de reconnaissance du pont exigeait un préfixe
\texttt{source \$SNAPSHOT\_FILE ... \&\&} obligatoire.

\vspace{0.2cm}
Mais un vrai appel observé en bout en bout \textbf{omet parfois entièrement}
cette ligne :
\begin{lstlisting}
bash -c -l "shopt -u extglob ... && eval '<commande>' ... && pwd -P >| ..."
(sans aucun `source` au debut)
\end{lstlisting}

\vspace{0.3cm}
\textbf{Correction} : le préfixe \texttt{source ...} est rendu
\textbf{optionnel} dans le motif. La condition exacte qui déclenche l'une ou
l'autre forme reste incomprise --- notée comme limite connue.
\end{frame}

% ============================================================
\section{Résultat actuel}

\begin{frame}[fragile]{Ce qui fonctionne maintenant}
Après les 5 corrections ci-dessus, l'outil \textbf{Bash} de Claude Code
fonctionne correctement à travers le pont ARIES, de bout en bout :
\begin{itemize}
  \item Démarrage / arrêt du harness --- \checkmark
  \item Authentification (\texttt{apiKeyHelper}) --- \checkmark
  \item Commandes Bash routées vers le sandbox --- \checkmark
  \item Plus aucune erreur \texttt{"exec request failed"} --- \checkmark
\end{itemize}

\vspace{0.3cm}
\texttt{status=succeeded}, cycle complet \texttt{Start}/\texttt{Run}/\texttt{Stop}
validé sur une vraie tâche (créer un fichier, le lire).
\end{frame}

\begin{frame}{Mais : une limite architecturale découverte par Claude Code lui-même}
Une fois Bash réparé, la réponse finale de l'agent a signalé :

\begin{quote}
\small ``\textit{the Bash tool is connected via SSH to a separate machine
(root@172.18.0.1, home dir /root) that doesn't share a filesystem with the one
my Write/Read tools use (/home/aries/workspace).}''
\end{quote}

\vspace{0.3cm}
\textbf{Exact} : les outils natifs \texttt{Write}/\texttt{Read}/\texttt{Edit}
de Claude Code manipulent le système de fichiers \textbf{directement}
(Node.js \texttt{fs}), sans jamais invoquer \texttt{bash} --- donc jamais
interceptés par le wrapper.

\vspace{0.3cm}
\textbf{Conséquence} : un fichier créé avec \texttt{Write} existe dans le
harness, jamais dans le sandbox --- invisible pour un vérificateur de
benchmark réel.
\end{frame}

\begin{frame}{État du code à date}
\begin{itemize}
  \item \texttt{pkg/harness/claudecode/} et \texttt{pkg/bridge/claudecodessh/} :
        écrits, compilent, testés contre de vraies infrastructures
  \item Outil Bash : \textbf{validé de bout en bout}
  \item Outils Write/Read/Edit : \textbf{non résolus} --- limite documentée,
        pas encore de correctif
  \item Documentation complète : \texttt{rapport\_integration\_claude\_code.md}
        (contexte, décisions de conception, tous les résultats empiriques
        détaillés ci-dessus)
\end{itemize}
\end{frame}

% ============================================================
\section{Pour la suite}

\begin{frame}{Questions pour discussion}
\begin{itemize}
  \item Résoudre Write/Read/Edit (serveur MCP personnalisé, ou montage
        réseau) fait-il partie du périmètre attendu du stage ?
  \item Continuer à approfondir Claude Code, ou consolider/documenter l'état
        actuel et passer à autre chose ?
\end{itemize}
\end{frame}

\begin{frame}{Merci}
\centering
\Large Questions ?

\vspace{1cm}
\normalsize
Documentation complète : \texttt{rapport\_integration\_claude\_code.md} \\
Code : \texttt{pkg/harness/claudecode/}, \texttt{pkg/bridge/claudecodessh/}
\end{frame}

\end{document}
